\chapter{DHCP: Client, Server and Relay}\label{ch:dhcp}
\bp{1.7}

\begin{outcomes}
By the end of this chapter you will be able to:
\begin{tight}
  \item \textbf{Describe} the four DHCPv4 messages and say which of them are
        broadcast and why that matters.
  \item \textbf{Configure} an IOS router in each of the three roles: client,
        server and relay agent.
  \item \textbf{Interpret} \cmd{show ip dhcp binding}, \cmd{show ip dhcp pool} and
        \cmd{show ip dhcp conflict}.
  \item \textbf{Troubleshoot} a host that has no address, has the wrong address, or
        has an address but no connectivity --- and tell those three apart before
        touching anything.
\end{tight}
\end{outcomes}

\begin{prereq}
\chref{ipv4} for addressing and subnet boundaries. \chref{vlans} and
\chref{trunking} for SVIs --- a DHCP fault is very often a VLAN fault wearing a
disguise. \chref{transport} for UDP and for broadcast versus unicast.
\end{prereq}

\begin{keyterms}
DORA $\bullet$ \cmd{DHCPDISCOVER} $\bullet$ lease $\bullet$ binding $\bullet$
relay agent $\bullet$ \cmd{ip helper-address} $\bullet$ \cmd{giaddr} $\bullet$
excluded address $\bullet$ APIPA $\bullet$ conflict
\end{keyterms}

\begin{examnote}[The verb is \emph{troubleshoot}, not \emph{configure}]
Topic~\tp{1.7} reads \emph{troubleshoot DHCPv4 client, server, and relay on IOS
devices}. It is worth reading twice. The exam is not primarily asking you to
build a DHCP pool from nothing --- it is asking you to be shown a pool, a relay
and a client that between them are not working, and to say which of the three is
at fault.

Note also \textbf{DHCPv4}. There is no DHCPv6 in this topic. IPv6 hosts in this
course get their addresses by SLAAC (\chref{ipv6}).

That is why this chapter spends more space on \emph{how to tell the three failure
classes apart} than on the pool syntax, which is four lines long.
\end{examnote}

\section{What DHCP is solving}

A host that has just booted knows its own MAC address and nothing else. It has no
IP address, so it cannot be the destination of a routed packet; it does not know
the address of a server, so it cannot ask one directly; and it does not know its
own subnet, so it cannot even work out what a sensible address would be.

Everything odd about DHCP follows from that opening position. The first message
has to be a broadcast, because there is no unicast address to use in either
direction. And broadcasts do not cross routers, which is the whole reason relay
agents exist.

\begin{definitionbox}[DHCPv4]
A client--server protocol that leases an IPv4 address and a set of configuration
options --- default gateway, DNS servers, domain name, lease time --- to a host
for a bounded period. The client runs on UDP port~68, the server on UDP port~67.
\end{definitionbox}

\subsection{DORA}

Four messages, in a fixed order, remembered as \term{DORA}.

\ccnafig{%
\begin{tikzpicture}[font=\scriptsize]
  \node[font=\scriptsize\bfseries, text=neutral!45!black] (cl) at (0,0)  {CLIENT};
  \node[font=\scriptsize\bfseries, text=neutral!45!black] (sv) at (10,0) {SERVER};
  \draw[neutral!50, line width=0.8pt] (0,-0.35) -- (0,-5.3);
  \draw[neutral!50, line width=0.8pt] (10,-0.35) -- (10,-5.3);

  \tikzset{
    bc/.style={-{Stealth[length=2.4mm]}, line width=1.1pt, dom4!80},
    uc/.style={-{Stealth[length=2.4mm]}, line width=1.1pt, dom1!80},
    lb/.style={font=\scriptsize\bfseries, above=1pt, midway},
    sb/.style={font=\tiny, text=neutral!85!black, below=1pt, midway}}

  \draw[bc] (0,-1.0)  -- node[lb, text=dom4!45!black] {1. DHCPDISCOVER}
                         node[sb] {broadcast, src 0.0.0.0, dst 255.255.255.255} (10,-1.0);
  \draw[bc] (10,-2.1) -- node[lb, text=dom4!45!black] {2. DHCPOFFER}
                         node[sb] {broadcast --- the client still has no address} (0,-2.1);
  \draw[bc] (0,-3.2)  -- node[lb, text=dom4!45!black] {3. DHCPREQUEST}
                         node[sb] {broadcast, so other servers withdraw their offers} (10,-3.2);
  \draw[uc] (10,-4.3) -- node[lb, text=dom1!45!black] {4. DHCPACK}
                         node[sb] {the lease is now a binding on the server} (0,-4.3);
\end{tikzpicture}}%
{DORA. Three of the four messages are broadcasts, which is why a router between
client and server stops the process dead unless it is configured to relay.}{dora}

Two details in that figure earn their place on the exam.

\begin{tight}
  \item \textbf{The REQUEST is broadcast, not unicast.} If two servers both
        offered, the request names the one the client accepted, and the other
        server sees it and returns its offered address to the pool. A unicast
        request would leave the second address reserved for nothing.
  \item \textbf{The client is not finished at the OFFER.} A host that receives an
        offer and then loses the server still has no address. ``It got an offer''
        is not the same as ``it is working''.
\end{tight}

\begin{conceptbox}[Renewal is not DORA]
At half the lease time (T1), a client that already has an address sends a
\emph{unicast} \cmd{DHCPREQUEST} straight to its server and expects an ACK. There
is no discover and no offer. This matters for troubleshooting: a fault that only
breaks broadcast forwarding --- a dead relay, a wrong VLAN --- can leave existing
hosts working perfectly for hours while every newly booted host fails. ``It works
for everyone except the machines that rebooted this morning'' is a relay fault
until proved otherwise.
\end{conceptbox}

\section{Role 1: IOS as a client}

A router interface can take its address from DHCP. This is normal on the WAN
interface of a small branch router.

\begin{config}{a router that leases its own WAN address}
interface GigabitEthernet0/0/0
 description WAN to ISP
 ip address dhcp
 no shutdown
\end{config}

\begin{verify}{R1}{show ip interface brief}
Interface              IP-Address      OK? Method Status                Protocol
GigabitEthernet0/0/0   203.0.113.24    YES DHCP   up                    up
GigabitEthernet0/0/1   192.168.10.1    YES manual up                    up
\end{verify}

The \cmd{Method} column is the whole point of that output. \cmd{DHCP} means the
address was leased; \cmd{manual} means somebody typed it. A WAN interface showing
\cmd{unassigned} with the line protocol up is a router that asked and got no
answer.

\begin{examnote}
A router that is a DHCP client normally also wants the default route the lease
carries. It installs one automatically as a static default with an administrative
distance of~254, so anything you configure yourself will win. Check with
\cmd{show ip route static} if a branch router is reaching the internet by a path
you did not expect.
\end{examnote}

\section{Role 2: IOS as a server}

Small sites often use the router itself rather than a dedicated server. Four
things must be right, and they must be right \emph{together}.

\begin{config}{a DHCP pool for VLAN 10}
! 1. Reserve the addresses that are already spoken for -- the gateway,
!    the switch management SVIs, the printers, the servers.
!    Do this FIRST. A pool created before the exclusions can hand out
!    the gateway address in the seconds between the two commands.
ip dhcp excluded-address 192.168.10.1 192.168.10.20
ip dhcp excluded-address 192.168.10.250 192.168.10.254
!
ip dhcp pool VLAN10-STAFF
 network 192.168.10.0 255.255.255.0
 default-router 192.168.10.1
 dns-server 192.168.20.10 192.168.20.11
 domain-name campus.example.edu
 lease 0 8 0
!
! lease is days hours minutes. "lease 0 8 0" is eight hours, which
! suits a teaching lab; "lease infinite" suits nothing.
\end{config}

\begin{alertbox}[The \cmd{network} statement must match the interface's subnet]
The router decides which pool to use by looking at the address of the interface
the request arrived on --- or, for a relayed request, at the \cmd{giaddr} field.
If VLAN~10's SVI is \cmd{192.168.10.1/24} and the pool says
\cmd{network 192.168.10.0 255.255.255.0}, they match and the pool is used.

If somebody later renumbers the SVI to \cmd{/25} and leaves the pool at
\cmd{/24}, the pool no longer matches, the router silently declines to serve, and
the symptom is a host with no address at all --- with a perfectly valid-looking
pool sitting in the configuration.
\end{alertbox}

\begin{verify}{R1}{show ip dhcp binding}
Bindings from all pools not associated with VRF:
IP address      Client-ID/              Lease expiration        Type      State
                Hardware address
192.168.10.21   0100.5079.6668.01       Aug 28 2026 06:14 AM    Automatic Active
192.168.10.22   0100.5079.6668.22       Aug 28 2026 06:19 AM    Automatic Active
\end{verify}

\begin{verify}{R1}{show ip dhcp pool}
Pool VLAN10-STAFF :
 Utilization mark (high/low)    : 100 / 0
 Subnet size (first/next)       : 0 / 0
 Total addresses                : 254
 Leased addresses               : 2
 Excluded addresses             : 25
 Pending event                  : none
 1 subnet is currently in the pool:
 Current index   IP address range                    Leased/Excluded/Total
 192.168.10.23   192.168.10.1     - 192.168.10.254    2  / 25 / 254
\end{verify}

\cmd{Leased addresses} against \cmd{Total addresses} is the number to read when
somebody reports that ``DHCP has stopped working'' at nine in the morning. A pool
that is full does not log an error to anybody who is watching; it simply stops
answering.

\section{Role 3: IOS as a relay agent}

The DISCOVER is a broadcast, and a router does not forward broadcasts. So when
the server is on a different subnet from the client --- which is the normal case
in any site with more than one VLAN --- the router on the client's subnet has to
pick the broadcast up and forward it as a unicast.

\begin{topology}
  \Host{h}{(0,0)}{PC1}
  \Sw{sw}{(2.6,0)}{SW1}
  \Rtr{r}{(6.0,0)}{R1}
  \Srv{s}{(9.6,0)}{DHCP\\10.0.0.5}
  \wire{h}{sw}{VLAN 10}
  \wire{sw}{r}{trunk}
  \wire{r}{s}{Gi0/1}
  \node[font=\tiny, text=dom4!45!black, align=center] at (2.6,-1.15)
    {\textbf{broadcast}\\stops here};
  \node[font=\tiny, text=dom1!45!black, align=center] at (7.8,-1.15)
    {\textbf{unicast}\\to 10.0.0.5};
\end{topology}

\begin{config}{relaying VLAN 10's requests to a central server}
interface Vlan10
 description Staff -- clients live here
 ip address 192.168.10.1 255.255.255.0
 ip helper-address 10.0.0.5
\end{config}

\begin{alertbox}[\cmd{ip helper-address} goes on the interface facing the
\emph{clients}]
This is the single most common relay mistake and it is worth stating on its own.
The command belongs on the SVI or interface in the subnet where the broadcasts
are --- the client side. Putting it on the interface facing the server does
nothing at all, because no DISCOVER ever arrives there.

If a site has six client VLANs, the command goes on six SVIs.
\end{alertbox}

\begin{conceptbox}[What \cmd{giaddr} does, and why the server needs it]
When the relay forwards the DISCOVER it fills in the \term{gateway IP address}
field, \cmd{giaddr}, with its own address on the interface that received the
broadcast --- \cmd{192.168.10.1} in the example above.

That single field does two jobs. It tells the server which subnet the client is
on, so the server knows which pool to draw from; and it gives the server a
unicast address to send the OFFER back to, so the relay can broadcast it onward
to the client.

The consequence for troubleshooting: if the central server has no pool for
\cmd{192.168.10.0/24}, the relay is working perfectly and the client still gets
nothing. Do not conclude ``the relay is broken'' from ``the client got no
address''.
\end{conceptbox}

\begin{examnote}[\cmd{ip helper-address} forwards more than DHCP]
By default it relays eight UDP services, not one --- among them TFTP~(69),
DNS~(53), NetBIOS name service~(137) and TACACS~(49) as well as DHCP~(67/68).
That is occasionally useful and occasionally a surprise, and it is a fair exam
question. Trim it with \cmd{ip forward-protocol udp <port>} if a site wants only
DHCP relayed.
\end{examnote}

\section{Troubleshooting: three symptoms, three places to look}

This is the part the blueprint actually asks for. Almost every DHCP fault
presents as one of three symptoms, and each points at a different part of the
path. Establish which one you have \emph{before} you start changing
configuration.

\begin{longtable}{@{}L{3.3cm}L{5.5cm}L{6.4cm}@{}}
\toprule
\rowcolor{primary!15}
\textbf{Symptom on the host} & \textbf{What it proves} & \textbf{Where to look} \\
\midrule
\endhead
\cmd{169.254.x.x}, or ``no address'' &
The host asked and heard nothing back. The failure is before the server ever
formed an opinion. &
Client VLAN and access port first (\chref{tshootl2}); then the relay ---
is \cmd{ip helper-address} on the client-side SVI? Then whether the SVI is up. \\
\rowcolor{lightbg}
An address, but from the wrong subnet &
Something answered, but from the wrong pool --- or a rogue server answered
first. &
\cmd{giaddr}: the client is in a different VLAN from the one you think. Or an
unauthorised server on the segment --- see \chref{l2security}. \\
The right address, but nothing works &
DHCP succeeded. This is not a DHCP fault at all. &
The \emph{options}, not the address: \cmd{default-router} and
\cmd{dns-server}. A host with a good address and a wrong gateway reaches its own
VLAN and nothing else. \\
\bottomrule
\end{longtable}

\begin{alertbox}[The third row is the one people waste an afternoon on]
``DHCP is broken'' is reported when a user cannot reach the internet. Run
\cmd{ipconfig /all} on the host. If it has an address in the right subnet, DHCP
has done its job perfectly and the problem is elsewhere --- routing, NAT, DNS or
an ACL. Check the gateway and DNS values the lease carried, then leave DHCP
alone.
\end{alertbox}

\subsection{Conflicts}

\begin{verify}{R1}{show ip dhcp conflict}
IP address       Detection method   Detection time
192.168.10.24    Ping               Aug 27 2026 09:12 AM
192.168.10.31    Gratuitous ARP     Aug 27 2026 09:41 AM
\end{verify}

A conflict means the router offered an address and then found somebody already
using it. \cmd{Ping} means the server tested the address before offering and got
a reply; \cmd{Gratuitous ARP} means a client took the address and immediately
heard another host claim it.

Either way the cause is the same: a statically addressed device inside the pool
range that nobody excluded. Conflicted addresses are removed from the pool and
stay out, so a site that never clears them slowly runs out of addresses. Fix the
exclusions first, then \cmd{clear ip dhcp conflict *}.

\begin{pitfall}
\begin{tight}
  \item \textbf{Pool created before the exclusions.} There is a real window in
        which the gateway address can be leased to a laptop. Exclude first.
  \item \textbf{\cmd{ip helper-address} on the server-side interface.} Does
        nothing. It goes on the client side.
  \item \textbf{Pool \cmd{network} not matching the SVI's subnet or mask.} Silent
        refusal to serve.
  \item \textbf{Assuming a working host proves DHCP works.} Existing leases renew
        by unicast and survive faults that stop every new host dead.
  \item \textbf{Forgetting the SVI must be up.} An SVI with no access port in its
        VLAN in the up state is down, so the pool it serves is unreachable ---
        a layer 2 fault presenting as a DHCP fault.
  \item \textbf{Reading \cmd{169.254.x.x} as a server problem.} It means nothing
        answered. Nine times in ten the fault is on the client's own switch port.
\end{tight}
\end{pitfall}

\begin{breakfix}[Every host in VLAN 30 gets 169.254 addresses. VLAN 10 and VLAN 20 are fine.]
The DHCP server is a central router, \cmd{R1}, reached through a distribution
switch. VLANs 10, 20 and 30 all have SVIs on \cmd{DSW1}, all three have pools on
\cmd{R1}, and only VLAN~30 fails. Static addressing inside VLAN~30 works, and a
statically addressed host in VLAN~30 can ping \cmd{R1}.

\begin{verify}{DSW1}{show run interface vlan 30}
interface Vlan30
 description Lab -- clients
 ip address 192.168.30.1 255.255.255.0
!
\end{verify}

\textbf{Diagnosis.} There is no \cmd{ip helper-address}. VLAN~10 and VLAN~20 have
one; VLAN~30 was added later and the line was missed. Everything else about
VLAN~30 is correct, which is exactly why static addressing works and why the
subnet is reachable --- routing is fine, layer~2 is fine, and only the broadcast
relay is absent.

The tell is that the fault is confined to \emph{one} VLAN. A server fault or a
routing fault would take out all three; a switch-port fault would take out one
host. One whole VLAN failing while its neighbours work points at something
configured per-SVI, and on this path there is only one such thing.

\textbf{Fix.}

\begin{config}{restoring the relay}
interface Vlan30
 ip helper-address 10.0.0.5
\end{config}

\textbf{Then check the second half.} Confirm \cmd{R1} actually has a pool whose
\cmd{network} statement matches \cmd{192.168.30.0/24}. A relay pointing at a
server with no matching pool produces the identical symptom, and fixing one half
of a two-part fault and declaring victory is how faults come back the next
morning.
\end{breakfix}

\begin{lab}{DHCP in all three roles, then broken deliberately}{Packet Tracer}
\textbf{Topology.} \cmd{SW1} carrying VLANs 10, 20 and 30; \cmd{DSW1} holding the
three SVIs; \cmd{R1} as the central server on \cmd{10.0.0.0/24}; two PCs per VLAN.

\begin{tasks}
  \item Address the SVIs and confirm inter-VLAN routing works with static host
        addresses before introducing DHCP at all. Do not skip this --- it
        separates the layer~3 work from the DHCP work.
  \item On \cmd{R1}, exclude the first twenty addresses of each subnet, then build
        three pools with gateway, DNS and an eight-hour lease.
  \item Add \cmd{ip helper-address} to VLAN~10's SVI only. Boot one PC in VLAN~10
        and one in VLAN~20. Record what each receives, and explain the difference
        in terms of \cmd{giaddr} before adding the second helper.
  \item Add the remaining helpers. Confirm all six PCs lease correctly and read
        \cmd{show ip dhcp binding} and \cmd{show ip dhcp pool} on \cmd{R1}.
  \item \textbf{Break 1.} Change VLAN~20's SVI mask to \cmd{/25}, leaving the pool
        at \cmd{/24}. Release and renew a VLAN~20 host. Explain the result.
  \item \textbf{Break 2.} Give a PC a static address inside a pool range, then
        force new leases until a conflict is detected. Read
        \cmd{show ip dhcp conflict}, fix the exclusion, and clear it.
  \item \textbf{Break 3.} Remove \cmd{default-router} from VLAN~30's pool. Renew a
        host. It will lease successfully and reach its own VLAN only. Write down
        the exact test that distinguishes this from a routing failure.
  \item \textbf{Extension.} Configure \cmd{R1}'s uplink as \cmd{ip address dhcp}
        against a second router acting as an ISP, and observe the default route
        that appears with distance 254.
\end{tasks}
\end{lab}

\begin{checkpoint}
\begin{tightnum}
  \item Which of the four DORA messages are broadcast, and what would break if
        the REQUEST were unicast?
  \item A host has \cmd{169.254.8.31}. What single fact does that establish, and
        what does it \emph{not} establish?
  \item On which interface does \cmd{ip helper-address} belong, and what does the
        relay write into \cmd{giaddr}?
  \item Existing hosts work; every rebooted host fails. What class of fault does
        that pattern point to, and why?
  \item A host has the correct address and gateway but cannot browse the web. Is
        this a DHCP fault? Justify the answer.
\end{tightnum}
\end{checkpoint}

\begin{chaptersummary}
\begin{tight}
  \item DORA: Discover, Offer, Request, Ack. The first three are broadcast, which
        is why routers block the process by default.
  \item Renewal at T1 is a \emph{unicast} request with no discover, so existing
        hosts survive faults that kill every new one.
  \item As a client: \cmd{ip address dhcp}, and the \cmd{Method} column of
        \cmd{show ip interface brief} proves it.
  \item As a server: exclude first, then a pool whose \cmd{network} statement
        matches the serving interface's subnet exactly.
  \item As a relay: \cmd{ip helper-address} on the \textbf{client-side}
        interface. It writes \cmd{giaddr}, which selects the pool and carries the
        reply back.
  \item Three symptoms, three places: no address means nothing answered; wrong
        subnet means the wrong pool or a rogue server; right address but no
        connectivity is not a DHCP fault.
  \item Conflicts come from static addresses inside the pool range. Fix the
        exclusion, then \cmd{clear ip dhcp conflict *}.
\end{tight}
\end{chaptersummary}

\begin{reviewq}
  \item \textbf{(MCQ)} A client receives \cmd{169.254.14.2}. What has happened?
        \begin{tight}
          \item[(a)] The server offered an address from the wrong pool
          \item[(b)] The client received no DHCP reply and self-assigned
          \item[(c)] The lease expired and was not renewed
          \item[(d)] The relay agent rewrote \cmd{giaddr} incorrectly
        \end{tight}
  \item \textbf{(MCQ)} Where must \cmd{ip helper-address} be configured?
        \begin{tight}
          \item[(a)] On the interface facing the DHCP server
          \item[(b)] On the SVI in the clients' own subnet
          \item[(c)] On every interface of the router
          \item[(d)] Under the \cmd{ip dhcp pool} for that subnet
        \end{tight}
  \item \textbf{(Short)} Explain why the DHCPREQUEST is broadcast rather than sent
        directly to the server that made the offer.
  \item \textbf{(Short)} A pool is configured, the SVI is up, the helper address
        is correct, and clients still receive nothing. Name two configuration
        errors that would produce exactly this and say how you would distinguish
        them.
  \item \textbf{(Applied)} Write the complete configuration for a router serving
        two VLANs, \cmd{10.1.10.0/24} and \cmd{10.1.20.0/24}, reserving the first
        ten addresses of each for infrastructure, with a four-hour lease and two
        DNS servers. Then state which single command you would add if the server
        were moved to a third subnet.
  \item \textbf{(Applied)} You are handed \cmd{show ip dhcp pool} output showing
        \cmd{Leased addresses: 251} of \cmd{Total addresses: 254}. Users report
        intermittent failures. Give the immediate mitigation and the permanent
        fix, and say which you would do first and why.
\end{reviewq}
